\section*{अध्याय \ref{ch:b3:planetary-interiors-surfaces} --- ग्रहीय अभ्यंतर और सतहें}

\begin{solution}{exo:b3:planetary-interiors-surfaces:1}
$\tfrac52C\bar\rho = 2.5\times0.35\times5429 = 4750$; $x^2 = (4750 -
3300)/(5429 - 3300) = 0.681$, $x = 0.825$, क्रोड $0.825\times2440 =
\qty{2010}{km}$; $\rho_c = 3300 + 2129/0.825^3 = \qty{7090}{kg.m^{-3}}$.
\end{solution}

\begin{solution}{exo:b3:planetary-interiors-surfaces:2}
$m = 3500\times\tfrac43\pi(25)^3 = \qty{2.29e8}{kg}$; $E = \tfrac12 m v^2 =
\tfrac12\times\num{2.29e8}\times(\num{2.0e4})^2 = \qty{4.6e16}{J}$, या
$\num{4.6e16}/\num{4.2e15} = \qty{11}{Mt}$। नियम से प्राप्त क्रेटर का व्यास
$20\times\qty{50}{m} = \qty{1}{km}$ होता है।
\end{solution}

\begin{solution}{exo:b3:planetary-interiors-surfaces:3}
$N_1 = 120/\num{100000} = \qty{1.2e-3}{km^{-2}}$, सापेक्ष त्रुटि $1/\sqrt{120}
= \qty{9}{\percent}$। रैखिक भाग पर, $T = N_1/c = \num{1.2e-3}/\num{4.875e-4}
= \qty{2.46}{Ga}$ (घातांकी पद, $\num{8e-7}$, नगण्य है), और
$\sigma_T = T/\sqrt{120} = \qty{0.22}{Ga}$.
\end{solution}

\begin{solution}{exo:b3:planetary-interiors-surfaces:4}
क्षेत्रफल $4\pi(\num{1.8215e6})^2 = \qty{4.17e13}{m^2}$; शक्ति $2.5\times
\num{4.17e13} = \qty{1.04e14}{W} = \qty{104}{TW}$, जो पृथ्वी की कुल शक्ति का 2.2 गुना है,
बारह गुना छोटे पृष्ठ से निकलती है, और पृथ्वी के फ़्लक्स से $2.5/0.092 = 27$ गुना
अधिक है।
\end{solution}

\begin{solution}{exo:b3:planetary-interiors-surfaces:5}
(क)~$\bar\rho = \num{6.417e23}/(\tfrac43\pi(\num{3.3895e6})^3) =
\qty{3934}{kg.m^{-3}}$; $g = GM/R^2 = \qty{3.73}{m.s^{-2}}$; $\rho gR =
\qty{4.97e10}{Pa}$। (ख)~$19\mu/(2\rho gR) = 19\times\num{6.7e10}/\num{9.94e10}
= 12.8$, $k_2 = 1.5/13.8 = 0.109$। (ग)~मापा गया मान
2.6 मानक विचलन से बड़ा है: मंगल एक ऐसे ठोस ग्रह से अधिक मान देता है जिसकी प्रावार दृढ़ता हो,
जैसा भीतर द्रव क्रोड होने पर होता।
\end{solution}

\begin{solution}{exo:b3:planetary-interiors-surfaces:6}
(क)~चंद्रमा: $4\times9.8/1.62 = \qty{24}{km}$ क्रिस्टलीय चट्टान में,
\qty{12}{km} अवसाद में; मंगल और बुध: $4\times9.8/3.7 =
\qty{10.6}{km}$ और \qty{5.3}{km}। चंद्रमा की \qty{20}{km} गहराई उसके दो पूर्वानुमानों
के बीच पड़ती है: नियम कुछ दसियों प्रतिशत तक सही है। (ख)~$Y =
\rho gD/3 = 2600\times9.8\times4000/3 = \qty{34}{MPa}$, विभंजित चट्टान की
सामर्थ्य है, न कि अक्षत क्रिस्टल की।
\end{solution}

\begin{solution}{exo:b3:planetary-interiors-surfaces:7}
मंगल: $\num{3.22e-14}(e^{6.93T} - 1) + \num{4.875e-4}\,T = \num{2.0e-3}$;
रैखिक अनुमान $4.1$~Ga बहुत अधिक पुराना है, क्योंकि
वहाँ घातांकी पद बड़ा है; द्विभाजन विधि से $T = \qty{3.34}{Ga}$ मिलता है (घातांकी पद
$\num{3.6e-4}$, रैखिक पद $\num{1.63e-3}$)। चंद्र फलन से प्राप्त
मान \qty{2.39}{Ga} है। प्रत्येक फलन एक ही पिंड के लिए अंशांकित होता है: संघट्टों की दर
और किसी दिए गए प्रक्षेप्य से बने क्रेटरों का आकार
चंद्रमा और मंगल के बीच भिन्न होते हैं।
\end{solution}

\begin{solution}{exo:b3:planetary-interiors-surfaces:8}
(क)~$3.551181/1.769138 = 2.0073$ और $7.154553/3.551181 = 2.0147$, दोनों
$2$ के निकट हैं। (ख)~$0.565247 - 0.844789 + 0.279542 = 0.000000$ प्रति दिन, $\num{1e-7}$
की सीमा तक: तीनों माध्य गतियाँ बद्ध हैं, एक ऐसा अनुनाद जो प्रत्येक संयोजन पर उपग्रहों की
सापेक्ष स्थितियाँ बनाए रखता है। (ग)~खिंचाव कक्षा के
उन्हीं बिंदुओं पर दोहराते हैं, इसलिए वे औसतन निरस्त होने के बजाय जुड़ते हैं;
वे उत्केंद्रता को उतनी ही तेज़ी से बढ़ाते हैं जितनी तेज़ी से ज्वार उसे क्षीण करते हैं।
\end{solution}

\begin{solution}{exo:b3:planetary-interiors-surfaces:9}
$n = 2\pi/(1.769138\times86400) = \qty{4.111e-5}{s^{-1}}$, $nR =
\qty{74.9}{m.s^{-1}}$, $(nR)^5/G = \qty{3.53e19}{W}$;
$\tfrac{21}{2}e^2(nR)^5/G = 10.5\times\num{1.6e-5}\times\num{3.53e19} =
\qty{5.9e15}{W}$ प्रति एकक $k_2/Q$। फिर $k_2/Q =
\num{1.04e14}/\num{5.9e15} = 0.018$.
\end{solution}

\begin{solution}{exo:b3:planetary-interiors-surfaces:10}
(क)~$x^2 = (\tfrac52C\bar\rho - \rho_m)/(\bar\rho - \rho_m)$ यह
$C$ में रैखिक है, जिसकी प्रवणता $\tfrac52\bar\rho/(\bar\rho - \rho_m)$ है। (ख)~प्रवणता
$2.5\times3340/52 = 161$ है: $x^2$ बदलता है $0.161$ से, $0.001$ के $C$ अंतर पर।
$C = 0.395$ के लिए, $x^2 = 0.036 + 0.161 = 0.197$, $x = 0.44$, $\rho_c =
\qty{3880}{kg.m^{-3}}$; $C = 0.393$ के लिए, $x^2 < 0$: बिल्कुल कोई क्रोड नहीं।
(ग)~$C = \tfrac25$ के निकट जड़त्व आघूर्ण एक छोटे क्रोड को मुश्किल से महसूस करता है:
इसका तीसरा सार्थक अंक तय करता है कि एक बड़ा हल्का क्रोड है या कोई भी नहीं।
भूकंपीय तरंगों, ज्वार और चुंबकीय क्षेत्र की आवश्यकता है।
\end{solution}

\begin{solution}{exo:b3:planetary-interiors-surfaces:11}
(क)~$\sigma_T = N_1/(10\,N_1')$: \qty{1.0}{Ga} पर, $N_1 \approx N_1' \approx
\num{4.875e-4}$, $\sigma_T = \qty{0.10}{Ga}$; \qty{3.8}{Ga} पर, $N_1 =
\num{1.07e-2}$ और $N_1' = \num{6.15e-2}$, $\sigma_T = \qty{0.017}{Ga}$,
छह गुना छोटा। (ख)~$N_1 = cT$, $N_1' = c$ मानकर $\sigma_T =
cT/(\sqrt n\,c) = T/\sqrt n$। (ग)~स्वयं कालक्रम फलन
अनिश्चित है: यह चंद्रमा के कुछ काल-निर्धारित स्थलों और चंद्रमा से मंगल तक के पैमाने के अंतरण पर टिका है,
ऐसी त्रुटियाँ जिन्हें कोई गिनती घटा नहीं सकती।
\end{solution}

\begin{solution}{exo:b3:planetary-interiors-surfaces:12}
(क)~अनुपात $(n_ER_E/n_IR_I)^5(e_E/e_I)^2 = (0.4982\times0.8569)^5
\times(0.009/0.004)^2 = 0.0142\times5.06 = 0.072$: लगभग
\qty{7.5e12}{W}, या $\num{7.5e12}/(4\pi(\num{1.5608e6})^2) =
\qty{0.24}{W.m^{-2}}$। (ख)~एक महासागर बर्फ़ की परत को
अंतःभाग से पृथक करता है: परत किसी ठोस चंद्रमा से कहीं अधिक मुड़ती है, और उसकी बर्फ़
आयो की गर्म चट्टान से भिन्न प्रकार से ऊर्जा का क्षय करती है, अतः $k_2$ और $Q$ दोनों बदलते हैं।
\end{solution}

\begin{solution}{pb:b3:planetary-interiors-surfaces:1}
\textbf{1.} $\bar\rho = \num{6.4169e23}/(\tfrac43\pi(\num{3.3895e6})^3) =
\qty{3934}{kg.m^{-3}}$.

\textbf{2.} $\tfrac52C\bar\rho = 3600$; $x^2 = (3600 - 3500)/(3934 - 3500) =
0.230$, $x = 0.479$, \qty{1620}{km} का क्रोड; $\rho_c = 3500 +
434/0.479^3 = \qty{7450}{kg.m^{-3}}$.

\textbf{3.} $x = 1830/3389.5 = 0.540$; $\rho_m = (3600 - 0.2915\times
3934)/(1 - 0.2915) = \qty{3460}{kg.m^{-3}}$, $\rho_c = 3462 +
472/0.1574 = \qty{6460}{kg.m^{-3}}$.

\textbf{4.} यह एक तिहाई कम है। मंगल के क्रोड में दाब
पृथ्वी के क्रोड की तुलना में बहुत कम है, अतः उसकी धातु कम संपीड़ित होती है; और इतना हल्का क्रोड
संकेत देता है कि लोहे के साथ कहीं अधिक हल्के तत्व मिले हुए हैं।

\textbf{5.} \omterm{def:b3:planetary-interiors-surfaces:love-number}{लव संख्या} (2003) इनसाइट (2021) से बहुत पहले सामने आ गई थी: उसने
पहले ही दिखा दिया था कि कम से कम बाह्य क्रोड द्रव है, और उसकी त्रिज्या
1520 और \qty{1840}{km} के बीच रखी, एक ऐसा परास जिसमें
भूकंपीय मान समाहित है।

\textbf{6.} $1 - (1 - 20/3389.5)^3 = \qty{1.8}{\percent}$ से $1 - (1 -
37/3389.5)^3 = \qty{3.2}{\percent}$ तक आयतन का, और द्रव्यमान का अंश उससे कम:
प्रावार से हल्का होने के कारण यह प्रभावी प्रावार घनत्व को थोड़ा घटाता है,
मॉडल की अनिश्चितता से छोटा एक प्रभाव है।

\textbf{7.} \qty{1}{km} से: 54; 1.41 से: 16; 2 से: 10; 2.83 से: 5;
4, 5.66, 8 और 11.3 से: 2; 16 से: 1.

\textbf{8.} $N = 54/80000 = \qty{6.75e-4}{km^{-2}}$, $\sigma = \sqrt{54}/80000
= \qty{0.92e-4}{km^{-2}}$ (\qty{14}{\percent}).

\textbf{9.} अपेक्षित $54\times0.355 = 19$ क्रेटर \qty{1.41}{km} से, और
$54\times0.147 = 7.9$ \qty{2}{km} से; गिने गए 16 और 10, एक
मानक विचलन के भीतर ($\sqrt{19} = 4.4$, $\sqrt{7.9} = 2.8$): संगत है।

\textbf{10.} नहीं: यह $0.072/\num{6.75e-4} = 107$ गुना कम है। गणना
समय दर्ज करती है।

\textbf{11.} $\log D$ में समान चौड़ाई के वर्ग परास को समान रूप से ढकते हैं और
तीव्र बंटन के वर्गों में तुलनीय संख्या में क्रेटर रखते हैं।

\textbf{12.} ये द्वितीयक क्रेटर हैं, जो एक बड़े प्राथमिक क्रेटर से
बाहर फेंके गए शिलाखंडों द्वारा बने हैं; ये अंतरिक्ष से आए स्वतंत्र संघट्ट नहीं हैं,
अतः ये ऐसे क्रेटर जोड़ देते जिन्हें कालक्रम फलन नहीं गिनता।

\textbf{13.} $T = \num{6.75e-4}/\num{4.875e-4} = \qty{1.385}{Ga}$; घातांकीय
पद, $\num{3.22e-14}(e^{9.6} - 1) = \num{4.7e-10}$, $N_1$ का
दस लाखवाँ भाग है।

\textbf{14.} $\sigma_T = T/\sqrt{54} = \qty{0.19}{Ga}$: $1.38 \pm 0.19$~Ga.

\textbf{15.} $\num{6.75e-4}/\num{8.38e-4} = \qty{0.81}{Ga}$, यह
लगभग आधा छोटा है।

\textbf{16.} $31 + 12 + 7 + 2 + 0 + 2 + 2 = 56$ क्रेटर; $N(\qty{2}{km}) =
56/120000 = \qty{4.67e-4}{km^{-2}}$; $N_1 = \num{4.67e-4}/0.1473 =
\qty{3.17e-3}{km^{-2}}$। $\num{3.22e-14}(e^{6.93T} - 1) +
\num{4.875e-4}\,T = \num{3.17e-3}$ पर द्विभाजन से $T = \qty{3.54}{Ga}$ प्राप्त होता है
(घातांकीय पद \num{1.45e-3}, रैखिक पद \num{1.73e-3})।

\textbf{17.} $N_1' = \num{3.22e-14}\times6.93\,e^{6.93\times3.54} +
\num{4.875e-4} = \num{1.05e-2}$; $\sigma_T = \num{3.17e-3}/(\sqrt{56}\times
\num{1.05e-2}) = \qty{0.040}{Ga}$.

\textbf{18.} मैदान A: \qty{14}{\percent}; मैदान B: \qty{1.1}{\percent}। कालक्रम के
तीव्र, प्रारंभिक भाग में क्रेटर घनत्व में बड़ा परिवर्तन
आयु में छोटे परिवर्तन का अर्थ है; रैखिक भाग पर आयु उतनी
ही अनिश्चित है जितनी गणना।

\textbf{19.} $120000\times\num{3.17e-3}\times(1 - 0.1473) = 324$; यह
$48 + 8 = 56$ क्रेटर ढोता है।

\textbf{20.} $N_1 = 56/(120000\times0.8527) = \qty{5.47e-4}{km^{-2}}$,
$T_r = \num{5.47e-4}/\num{4.875e-4} = \qty{1.12}{Ga}$.

\textbf{21.} $\sigma = 1.12/\sqrt{56} = \qty{0.15}{Ga}$.

\textbf{22.} \qty{2}{km} का क्रेटर \qty{400}{m} गहरा है: लावा कम से कम
उतना मोटा रहा होगा। बड़े क्रेटर अधिक गहरे हैं और उनके किनारे ऊँचे खड़े हैं;
उनमें से कुछ भर गए थे और अब भी दिखते हैं।

\textbf{23.} $1.38 \pm 0.19$ और $1.12 \pm 0.15$~Ga में \qty{0.26}{Ga} का अंतर है,
लगभग एक संयुक्त मानक विचलन ($\sqrt{0.19^2 + 0.15^2} = 0.24$) के बराबर है:
गणनाएँ एक ही घटना की अनुमति देती हैं, यद्यपि वे इसे सिद्ध नहीं करतीं।

\textbf{24.} युवा लावा मैदान A की मॉडल आयु
\textbf{1.4~अरब~वर्ष} ($1.38 \pm 0.19$~Ga) है।
\end{solution}
